\subsection{window-\texorpdfstring{$6$}{6} 的 visible Cartan 完成、syzygy Green 逆与模 \texorpdfstring{$2$}{2} 支撑间隙}\label{subsec:conclusion-window6-cartan-syzygy-green-isotropy-codegap}
沿用正文关于 window-\(6\) 可见向量族 \(\{v_w\}_{w\in X_6}\subset \ZZ^3\)、visible Cartan quotient
\[
\Pi_6:\ZZ^{X_6}\to \ZZ^3,
\qquad
\Pi_6(\delta_w)=v_w,
\]
以及 cyclic 与 full visible 矩阵
\[
A:=A_6^{\mathrm{cyc}}=[v_w]_{w\in X_6^{\mathrm{cyc}}}\in M_{3\times 18}(\ZZ),
\qquad
A_6:=[v_w]_{w\in X_6}\in M_{3\times 21}(\ZZ)
\]
的记号。前述链条已经分别给出
\[
AA^{\mathsf T}=10I_3,
\qquad
A_6A_6^{\mathsf T}=11I_3,
\]
syzygy 格 \(\Lambda_6^{\mathrm{syz}}=\ker \Pi_6\) 的判别式
\[
\operatorname{disc}(\Lambda_6^{\mathrm{syz}})=11^3,
\]
visible quartic counterterm 的唯一性、sextic 各向同性的绝对障碍，以及 Cartan kernel code / 对偶码的参数
\[
C_6^\perp\ \text{为 }[21,3,9],
\qquad
C_6\ \text{为 }[21,18,2].
\]
把这些对象进一步并置后，可把 window-\(6\) 的 rank-\(3\) visible Cartan 完成压缩为一条更紧的整数--实分析链：syzygy Green 算子是一个秩 \(3\) 偏差于恒等的显式对象；visible Cartan 荷的最小连续 lifting 被 \(1/11\) 的二次作用量严格锁定；syzygy 判别式、visible 协方差体积与全部三列子式平方预算共同塌缩到同一个 \(11^3\) 全息常数；quartic 阶仍存在唯一的精确 isotropization 界面，而 sextic 阶已彻底脱离径向补偿；模 \(2\) 阴影则把 Cartan 可见性与 syzygy 中性严格分裂到 \(9\)-局域与 \(2\)-局域两个支撑阈值。

\begin{theorem}[syzygy Green 逆与 cyclic Cartan 投影的显式闭式]\label{thm:conclusion-window6-syzygy-green-inverse-closed-form}
\leanverified{paper\_conclusion\_window6\_syzygy\_green\_inverse\_closed\_form}
记
\[
R_6:=\binom{I_{18}}{-A}\in M_{21\times 18}(\ZZ),
\qquad
K_6:=R_6^{\mathsf T}R_6=I_{18}+A^{\mathsf T}A.
\]
则有严格恒等式
\[
K_6^{-1}=I_{18}-\frac{1}{11}A^{\mathsf T}A.
\]
并且 cyclic 系数空间 \(\RR^{18}\) 上存在显式正交投影
\[
P_{\mathrm{Cartan}}=\frac{1}{10}A^{\mathsf T}A,
\qquad
P_{\mathrm{ker}}=I_{18}-\frac{1}{10}A^{\mathsf T}A,
\]
满足
\[
\operatorname{im}(P_{\mathrm{Cartan}})=\operatorname{im}(A^{\mathsf T}),
\qquad
\operatorname{im}(P_{\mathrm{ker}})=\ker A.
\]
特别地，window-\(6\) 的 cyclic Cartan 耦合只沿三条方向发生非平凡压缩，其余十五条方向在 syzygy Green 算子下保持单位刚性。
\end{theorem}
\begin{proof}
由
\[
AA^{\mathsf T}=10I_3
\]
立刻得到
\[
(A^{\mathsf T}A)^2=A^{\mathsf T}(AA^{\mathsf T})A=10A^{\mathsf T}A.
\]
因此
\[
\Bigl(I_{18}-\frac{1}{11}A^{\mathsf T}A\Bigr)\Bigl(I_{18}+A^{\mathsf T}A\Bigr)
=
I_{18}+\frac{10}{11}A^{\mathsf T}A-\frac{1}{11}(A^{\mathsf T}A)^2
=
I_{18},
\]
从而
\[
K_6^{-1}=(I_{18}+A^{\mathsf T}A)^{-1}=I_{18}-\frac{1}{11}A^{\mathsf T}A.
\]

再看
\[
\Bigl(\frac{1}{10}A^{\mathsf T}A\Bigr)^2
=
\frac{1}{100}(A^{\mathsf T}A)^2
=
\frac{1}{10}A^{\mathsf T}A,
\]
故 \(P_{\mathrm{Cartan}}\) 为投影。其像空间显然包含于 \(\operatorname{im}(A^{\mathsf T})\)；反之，对任意 \(u\in\RR^3\)，有
\[
P_{\mathrm{Cartan}}(A^{\mathsf T}u)
=
\frac{1}{10}A^{\mathsf T}A(A^{\mathsf T}u)
=
\frac{1}{10}A^{\mathsf T}(AA^{\mathsf T})u
=
A^{\mathsf T}u,
\]
故 \(\operatorname{im}(P_{\mathrm{Cartan}})=\operatorname{im}(A^{\mathsf T})\)。于是
\[
P_{\mathrm{ker}}=I_{18}-P_{\mathrm{Cartan}}
\]
也是投影，且其像空间恰为 \(\ker A\)。证毕。
\end{proof}

\begin{theorem}[visible Cartan 荷的 \texorpdfstring{$1/11$}{1/11} 极小提升律]\label{thm:conclusion-window6-cartan-charge-one-over-eleven-minimal-lift}
把 \(\Pi_6\) 线性延拓为
\[
\Pi_{6,\RR}:\RR^{X_6}\to \RR^3,
\qquad
\Pi_{6,\RR}(z)=A_6z.
\]
则对任意 \(u\in\RR^3\)，仿射纤维
\[
\mathcal L(u):=\{z\in\RR^{X_6}:A_6z=u\}
\]
中存在唯一极小 Euclid 范数解
\[
z_u^\star=\frac{1}{11}A_6^{\mathsf T}u,
\]
并且其最小作用量满足严格闭式
\[
\min_{z\in\mathcal L(u)}|z|_2^2=\frac{1}{11}|u|_2^2.
\]
因而 \(\Pi_{6,\RR}\) 的 Moore--Penrose 逆恰为
\[
\Pi_{6,\RR}^+=\frac{1}{11}A_6^{\mathsf T}.
\]
\end{theorem}
\begin{proof}
由
\[
A_6A_6^{\mathsf T}=11I_3
\]
可知 \(A_6\) 满行秩。于是 Moore--Penrose 伪逆的标准公式给出
\[
\Pi_{6,\RR}^+=A_6^{\mathsf T}(A_6A_6^{\mathsf T})^{-1}
=\frac{1}{11}A_6^{\mathsf T}.
\]
因此 \(\mathcal L(u)\) 中的唯一最小范数元为
\[
z_u^\star=\Pi_{6,\RR}^+u=\frac{1}{11}A_6^{\mathsf T}u.
\]
其范数由
\[
|z_u^\star|_2^2
=
u^{\mathsf T}(A_6A_6^{\mathsf T})^{-1}u
=
\frac{1}{11}|u|_2^2
\]
给出。唯一性来自满行秩情形下伪逆极小性的标准结论。证毕。
\end{proof}

\begin{corollary}[visible Cartan 荷的正交作用量分解]\label{cor:conclusion-window6-cartan-charge-orthogonal-energy-splitting}
对任意 \(z\in\RR^{X_6}\)，都有正交分解
\[
|z|_2^2
=
\bigl|z-z_{\Pi_{6,\RR}(z)}^\star\bigr|_2^2
+\frac{1}{11}\bigl|\Pi_{6,\RR}(z)\bigr|_2^2.
\]
故 visible Cartan 荷的连续审计代价被精确锁定为系数 \(1/11\) 的二次型。
\end{corollary}
\begin{proof}
由定理 \ref{thm:conclusion-window6-cartan-charge-one-over-eleven-minimal-lift}，\(z_{\Pi_{6,\RR}(z)}^\star\) 是纤维
\[
\mathcal L(\Pi_{6,\RR}(z))
\]
中的唯一最小范数元。最小范数分解定理给出
\[
|z|_2^2=\bigl|z-z_{\Pi_{6,\RR}(z)}^\star\bigr|_2^2+|z_{\Pi_{6,\RR}(z)}^\star|_2^2.
\]
再代入
\[
|z_{\Pi_{6,\RR}(z)}^\star|_2^2=\frac{1}{11}\bigl|\Pi_{6,\RR}(z)\bigr|_2^2
\]
即得所述公式。证毕。
\end{proof}

\begin{theorem}[syzygy 判别式、visible 协方差体积与三列子式预算的全息恒等式]\label{thm:conclusion-window6-syzygy-discriminant-volume-budget-holography}
对 visible 矩阵 \(A_6\) 与 syzygy 格 \(\Lambda_6^{\mathrm{syz}}\)，成立三重恒等式
\[
\operatorname{disc}(\Lambda_6^{\mathrm{syz}})
=
\det K_6
=
\det(A_6A_6^{\mathsf T})
=
\sum_{\substack{I\subset X_6\\ |I|=3}}\det(A_6[I])^2
=
11^3.
\]
其中 \(A_6[I]\) 表示取列集 \(I\) 所得的 \(3\times 3\) 子矩阵。
\end{theorem}
\begin{proof}
由定理 \ref{thm:window6-syzygy-gram-spectrum-discriminant}，
\[
\det K_6=\operatorname{disc}(\Lambda_6^{\mathrm{syz}})=11^3.
\]
另一方面，由
\[
A_6A_6^{\mathsf T}=11I_3
\]
得到
\[
\det(A_6A_6^{\mathsf T})=11^3.
\]
再由 Cauchy--Binet 公式，
\[
\det(A_6A_6^{\mathsf T})
=
\sum_{\substack{I\subset X_6\\ |I|=3}}\det(A_6[I])^2.
\]
合并上述等式即得结论。证毕。
\end{proof}

\begin{theorem}[quartic 可补偿而 sextic 不可补偿的各向同性分岔]\label{thm:conclusion-window6-quartic-sextic-isotropy-bifurcation}
定义 shell-constant quartic 空间
\[
\mathscr Q_4:=
\left\{
\mathcal Q_{a_2,a_3,a_4,b}(x):
a_2,a_3,a_4,b\in\RR
\right\},
\]
其中
\[
\mathcal Q_{a_2,a_3,a_4,b}(x)
:=
a_2\sum_{\beta\in \Phi_{6,2}^{\mathrm{cyc}}}\langle x,\beta\rangle^4
+a_3\sum_{\beta\in \Phi_{6,3}^{\mathrm{cyc}}}\langle x,\beta\rangle^4
+a_4\sum_{\beta\in \Phi_{6,4}^{\mathrm{cyc}}}\langle x,\beta\rangle^4
+b\sum_{i=1}^{3}x_i^4.
\]
又定义 visible sextic 空间
\[
\mathscr Q_6:=
\left\{
P_c(x):=\sum_{w\in X_6}c_w\langle x,v_w\rangle^6:\ c_w\in\RR
\right\}.
\]
则有
\[
\dim\bigl(\mathscr Q_4\cap \RR\cdot |x|^4\bigr)=1,
\qquad
\mathscr Q_6\cap \RR\cdot |x|^6=\{0\}.
\]
特别地，quartic 阶存在且仅存在一条各向同性补偿方向
\[
\sum_{w\in X_6^{\mathrm{cyc}}}\langle x,v_w\rangle^4
+2\sum_{i=1}^3x_i^4
=
12|x|^4,
\]
而 sextic 阶已发生不可逆的各向异性崩塌。
\end{theorem}
\begin{proof}
第一条正是定理 \ref{thm:window6-visible-quartic-shell-isotropization} 的直接重述：若
\[
\mathcal Q_{a_2,a_3,a_4,b}(x)=c|x|^4,
\]
则必且仅必
\[
a_2=a_3=a_4=:a,
\qquad
b=2a,
\qquad
c=12a.
\]
故交空间一维。

第二条由命题 \ref{prop:window6-visible-degree6-isotropy-obstruction} 给出。该命题证明：任意
\[
P_c(x)=\sum_{w\in X_6}c_w\langle x,v_w\rangle^6
\]
中，单项式 \(x_1^2x_2^2x_3^2\) 的系数恒为零；而 \(|x|^6\) 在该单项式上的系数等于 \(6\)。于是若 \(P_c=\lambda |x|^6\)，只能有 \(\lambda=0\)。证毕。
\end{proof}

\begin{theorem}[Cartan 可见性与 syzygy 中性的 \texorpdfstring{$9/2$}{9/2} 支撑间隙]\label{thm:conclusion-window6-cartan-syzygy-nine-two-support-gap}
令
\[
C_6^\perp=\operatorname{rowspan}_{\FF_2}(\bar A_6)\subset \FF_2^{21},
\qquad
C_6=\ker \bar A_6.
\]
则存在严格支撑间隙
\[
\operatorname{wt}(y)\ge 9
\qquad (0\neq y\in C_6^\perp),
\]
以及
\[
\exists\,0\neq z\in C_6:\ \operatorname{wt}(z)=2.
\]
因此：
\begin{itemize}
  \item 任一非平凡的 mod-\(2\) Cartan 见证都至少需要 \(9\) 个可见坐标；
  \item 非平凡的 syzygy 中性扰动已经可以在 \(2\) 个坐标上发生。
\end{itemize}
\end{theorem}
\begin{proof}
由定理 \ref{thm:window6-cartan-kernel-code-parameters}，
\[
C_6^\perp\ \text{是 }[21,3,9]\text{ 码},
\qquad
C_6\ \text{是 }[21,18,2]\text{ 码}.
\]
因此 \(C_6^\perp\) 的最小非零重量恰为 \(9\)，而 \(C_6\) 中存在重量为 \(2\) 的非零码字。证毕。
\end{proof}

\begin{corollary}[局部审计盲区]\label{cor:conclusion-window6-local-audit-blind-zone-for-cartan}
任意仅访问至多 \(8\) 个 visible 坐标的模 \(2\) 线性审计，都不可能检测到非平凡 Cartan 奇偶荷；反之，syzygy 中性误差却可在两点局部发生。于是 window-\(6\) 中存在严格的全局/局部分裂：
\[
\text{Cartan 可见性是 \(9\)-局域的，syzygy 中性却是 \(2\)-局域的。}
\]
\end{corollary}
\begin{proof}
若某个模 \(2\) 线性审计仅访问至多 \(8\) 个坐标，则其所有可见证向量的支撑都不超过 \(8\)。由定理 \ref{thm:conclusion-window6-cartan-syzygy-nine-two-support-gap}，任何非零 \(C_6^\perp\) 码字的重量都至少为 \(9\)，故这类审计不可能击中非平凡 Cartan 见证。另一方面，同一定理已给出重量为 \(2\) 的非零 \(C_6\) 码字存在性，故 syzygy 中性扰动确可局部发生。证毕。
\end{proof}

\begin{theorem}[极小 Cartan 见证的唯一性]\label{thm:conclusion-window6-unique-minimal-cartan-witness}
\(C_6^\perp\) 的非零码字按 Hamming 重量呈现刚性分层：
\[
W_{C_6^\perp}(z)-1=z^9+3z^{11}+3z^{14}.
\]
因此：
\begin{enumerate}
  \item 存在且仅存在一个最小支撑的非零 Cartan 见证，其重量恰为 \(9\)；
  \item 恰有三个非零见证重量为 \(11\)；
  \item 恰有三个非零见证重量为 \(14\)。
\end{enumerate}
特别地，全部七个非平凡 Cartan 奇偶见证并非任意散布，而是被严格锁定为 \(1+3+3\) 的重量分层。
\end{theorem}
\begin{proof}
由定理 \ref{thm:window6-cartan-kernel-code-parameters}，
\[
W_{C_6^\perp}(z)=1+z^9+3z^{11}+3z^{14}.
\]
于是非零码字总数为
\[
1+3+3=7=2^3-1,
\]
与 \(\dim C_6^\perp=3\) 一致。各重量层的基数恰为重量枚举多项式中相应系数。尤其重量 \(9\) 的系数等于 \(1\)，故最小支撑非零见证唯一。证毕。
\end{proof}

\begin{conjecture}[rank-\texorpdfstring{$r$}{r} visible Cartan 完成的 \texorpdfstring{$(c+1)^r$}{(c+1)^r} 全息刚性]\label{conj:conclusion-window-visible-cartan-rank-r-holographic-rigidity}
设更高窗情形中存在与 window-\(6\) 同型的 visible Cartan 数据：
\begin{enumerate}
  \item 一个分裂短正合列
  \[
  0\to \Lambda_m^{\mathrm{syz}}\to \ZZ^{N_m}\xrightarrow{\ \Pi_m\ }\ZZ^r\to 0;
  \]
  \item 一个 cyclic 矩阵 \(A_m^{\mathrm{cyc}}\in M_{r\times (N_m-r)}(\ZZ)\) 满足
  \[
  A_m^{\mathrm{cyc}}(A_m^{\mathrm{cyc}})^{\mathsf T}=c_mI_r;
  \]
  \item 一个 full visible 矩阵 \(A_m\in M_{r\times N_m}(\ZZ)\) 满足
  \[
  A_mA_m^{\mathsf T}=(c_m+1)I_r.
  \]
\end{enumerate}
则应当同时成立
\[
\operatorname{disc}(\Lambda_m^{\mathrm{syz}})=(c_m+1)^r,
\qquad
\Pi_{m,\RR}^+=\frac{1}{c_m+1}A_m^{\mathsf T},
\qquad
\sum_{|I|=r}\det(A_m[I])^2=(c_m+1)^r,
\]
并且其模 \(2\) visible Cartan 对偶码的非零重量分层仅由 boundary \(2\)-torsion 的轨道结构决定。window-\(6\) 正对应于
\[
r=3,
\qquad
c_m=10,
\qquad
(c_m+1)^r=11^3.
\]
\end{conjecture}

\paragraph{统一读法}
由此，window-\(6\) 的 rank-\(3\) visible Cartan 完成已不再只是某组彼此并列的可见模态。更精确地说，同一个参数族
\[
10,\qquad 11,\qquad 11^3,\qquad [21,3,9],\qquad [21,18,2]
\]
同时锁定了五层结构：cyclic Cartan 压缩只沿三条方向发生，syzygy Green 逆因此成为秩 \(3\) 偏差于恒等的显式对象；visible Cartan 荷的最小连续 lifting 被 \(1/11\) 的二次作用量严格控制；syzygy 判别式、visible 协方差体积与全部三列子式平方预算又共同塌缩到同一个全息常数 \(11^3\)；可见支持在 quartic 阶仍保留唯一的精确 isotropization 界面，而 sextic 阶已绝对脱离径向补偿；最后，mod-\(2\) 阴影迫使非平凡 Cartan 可见性至少在 \(9\) 个坐标上全局展开，而 syzygy 中性扰动却可以在 \(2\) 点局部发生，并且最小 Cartan 见证还是唯一的。于是，整数 Cartan 商空间、实二次最小作用量、三体体积预算、高阶各向同性分岔与模 \(2\) 支撑阈值便被压缩到同一条刚性主线上。

\endinput
